\chapter{The Aufhebung Decomposition}\label{ch:aufhebung}

\section{Why decompose composition?}\label{sec:why-decompose}

The bare monoidal datum $(\Ideas,\compose,\unit)$ tells us that ideas
combine, but it does not tell us \emph{how} they combine. When two
ideas $a$ and $b$ are placed in compositional relation, several things
happen at once: some structural features carry through unchanged from
either factor, some are explicitly cancelled or inverted, and some new
content---irreducible to either factor---comes into being. The history
of philosophy has a name for this triple movement; in
\citet{Hegel1812} it is \emph{Aufhebung}, the dialectical operation
that simultaneously preserves, negates, and elevates. The aim of this
chapter is to convert that informal triple into a piece of
mathematics, and to verify the conversion against seven historical
positions for which a precise notation has long been wanted but never
supplied.

The motivation is best seen from a failure mode of the bare monoid.
Consider an ideational lineage in which a \emph{thesis} $a$ is followed
by an \emph{antithesis} $\bar a$, yielding a composite $a\compose\bar
a$. Naively one might want $a\compose\bar a=\unit$: the thesis and its
negation cancel. But experience belies this. The Hegelian dialectic
denies it explicitly: the new product is not nothing, it is a
\emph{higher} idea $a'$ that retains a memory of $a$ and incorporates
the negation as a determinate moment. Likewise, in literary tradition,
a parody $\bar a$ of a canonical text $a$ does not annul $a$; the
parody's joint reading produces a new genre. In cognitive science,
recalling a memory $a$ under a transforming schema $\bar a$ does not
delete $a$; it produces a reconsolidated trace $a'$ in which the
original is preserved but altered \citep{Adorno1966}. The pattern is
the same in each case: the composition of two ideas yields neither
their union nor their cancellation but a third object whose anatomy
must be told.

What is missing from the bare monoid is a way to talk about the
\emph{anatomy} of $a\compose b$: which part of the composite is
attributable to ``what carried through,'' which part to ``what was
struck out,'' and which part to ``what was newly produced.'' Chapter 1
of this monograph already supplied the global $\sigma+\nu+\eta$
decomposition theorem in the resonance setting (Theorem 1.3,
``Aufhebung Decomposition''). The present chapter does something
different: it studies the decomposition as a structure on the carrier
$\Ideas$ itself, axiomatized minimally and analyzed exhaustively. We
formalize the \emph{intra-idea} version: every idea $a$ has a
preserved part $\sigma(a)$, a negated part $\nu(a)$, and the synthetic
residue (the elevation $\eta$) is recovered as the manner in which
they compose to give $a$. The reader will note that this is a
deliberate sharpening of the Hegelian text, which speaks of
\emph{moments} rather than \emph{functions}. Following
\citet{Brandom1994}'s reading, we treat $\sigma$ and $\nu$ as a pair
of commuting projectors on the conceptual lattice, and adopt the
strong axiom that $\sigma(a)\compose\nu(a)=a$ exactly. The cost of
this sharpening is some loss of generality at the philosophical level;
the gain is a clean algebraic theory whose roughly one hundred
elementary lemmas are all decidable from five axioms.

It is worth pausing on why the unary version of the decomposition
deserves a chapter of its own, given that the binary version is the
one that figures in the SHARED BRIEF's foundational nine theorems.
The reason is that the unary version is where the \emph{philosophical}
content lives. The binary $(\sigma,\nu,\eta)$ is a triple of
\emph{operations} on pairs; one cannot read off from it what it
means for an individual idea to be ``preserved'' or ``negated.'' The
unary projectors of \cref{def:aufheb-struct} below send each idea to
its preserved and negated \emph{parts}, and so make first-order
predicates---``$a$ is preserved,'' ``$a$ is sedimented''---available
for direct definition. These are the predicates the zoo entries
trade in. Hegel does not ask whether a binary operation is
$\sigma$-like; he asks whether an idea is \emph{aufgehoben}. The
unary structure formalizes that question; the binary structure
formalizes its consequences.

The split also resolves a tension between ``negation as cancellation''
and ``negation as transformation.'' The cross-annihilation axioms
$\sigma\circ\nu=\unit$ and $\nu\circ\sigma=\unit$ encode formally
that the preserved part of a negation is trivial, and the negated
part of a preservation is trivial---each projector recognizes the
\emph{other's} output as ``empty'' from its own point of view. This is
exactly Hegel's criterion of \emph{determinate} negation: a negation
preserves nothing of what it negates, and a preservation negates
nothing of what it preserves; what binds them is composition. The
chapter develops this minimal axiomatic in
\S\ref{sec:aufheb-defs}--\S\ref{sec:aufheb-iter}, and then exhibits
seven historical positions in the zoo of \S\ref{sec:aufheb-zoo} that
are made precise (or sharpened, or refuted) by it.

There is a methodological remark to enter before we proceed. The
Aufhebung structure of \cref{def:aufheb-struct} below is, in
mathematical taste, \emph{biased}. It privileges a ``preserved
first, negated second'' reading of composition: the decomposition
$a=\sigma(a)\compose\nu(a)$ writes the preserved part on the left,
and the cross-annihilation axioms send the projector outputs of one
side to the identity under the other side's projector. This bias is
not arbitrary; it tracks the order in which the dialectical moments
appear in Hegel's exposition (\emph{Sein} before \emph{Nichts},
position before negation, sediment before transformation). One could
develop a left--right symmetric variant in which both
$\sigma(a)\compose\nu(a)=a$ and $\nu(a)\compose\sigma(a)=a$ are
imposed, but the resulting structure forces composition to be
commutative on the projectors' joint image, which empirically fails
in the application domains of interest \citep[cf.][on the asymmetry
of metaphor]{LakoffJohnson1980}. We therefore retain the asymmetry
as a feature.

A second remark concerns scope. The decomposition we axiomatize is
\emph{unary}: it sends an idea to its preserved and negated parts.
The bilinear $\sigma(a,b),\nu(a,b),\eta(a,b)$ of the SHARED BRIEF is
the binary version, defined on the resonance-extended free vector
space in \cref{ch:foundations}. The two are related by restriction:
the unary $\sigma(a)$ here is the diagonal $\sigma(a,a)$ there, with
the same identification for $\nu$, and the elevation $\eta$ is
recovered as the residue of the binary product.  The unary version
is the right setting for the algebraic theorems of this chapter; the
binary version is the right setting for the homological computations
of \cref{ch:foundations}. We consistently use the unary $\sigma,\nu$
here and reserve the binary forms for cross-references.

A third remark concerns terminology. The English word ``sublation''
has been used by translators of Hegel since at least the nineteenth
century to render \emph{Aufhebung}; we keep it as a noun and use
\emph{sublated} for the predicate. The cognate ``sediment'' is
borrowed from sociology of knowledge \citep{Lovejoy1936}; an idea is
\emph{sedimented} when it has nothing left to negate, i.e.\ when its
$\nu$-component is the identity. The four predicates of
\cref{def:aufheb-preds}---preserved, negated, sedimented, sublated---
should be heard as a cross-classification of $a$ by the values of
$\sigma$ and $\nu$. Of the sixteen logically possible
combinations, only a small number are populated in any concrete
algebra; \cref{thm:3-3} below shows that several reduce to one
another.

\section{Definitions}\label{sec:aufheb-defs}

We work throughout with a fixed bare idea algebra
$(\Ideas,\compose,\unit)$ in the sense of
\cref{ch:foundations}: a set with an associative composition and a
two-sided identity. The Aufhebung structure is an additional layer.

\begin{definition}[Aufhebung structure]\label{def:aufheb-struct}
An \emph{Aufhebung structure} on $(\Ideas,\compose,\unit)$ consists
of two endofunctions
\[
\sigma:\Ideas\to\Ideas,\qquad \nu:\Ideas\to\Ideas
\]
satisfying the following five axioms for all $a\in\Ideas$:
\begin{enumerate}
\item[(A1)] \emph{Decomposition.} $\sigma(a)\compose\nu(a)=a$.
\item[(A2)] \emph{$\sigma$-idempotence.} $\sigma(\sigma(a))=\sigma(a)$.
\item[(A3)] \emph{$\nu$-idempotence.} $\nu(\nu(a))=\nu(a)$.
\item[(A4)] \emph{Identity preservation.} $\sigma(\unit)=\unit$ and $\nu(\unit)=\unit$.
\item[(A5)] \emph{Cross annihilation.}
$\sigma(\nu(a))=\unit$ and $\nu(\sigma(a))=\unit$.
\end{enumerate}
We call $\sigma(a)$ the \emph{preserved} (or \emph{sedimented})
component, $\nu(a)$ the \emph{negated} (or \emph{transformed})
component, and refer to (A1) as the \emph{Aufhebung law}.
\end{definition}

The synthesis or elevation $\eta(a)$ is implicit in (A1): given the
preserved part and the negated part, $a$ is reconstructed by
composition, and the \emph{way} they compose---the curvature of the
binary product $\compose$ at the pair $(\sigma(a),\nu(a))$---is what
the global theorem of \cref{ch:foundations} packages as $\eta$. From
the present, intra-idea, vantage, $\eta$ is not a separate function
but the residual datum carried by the monoid operation itself.

\begin{definition}[Aufheb pair]\label{def:aufheb-pair}
For $a\in\Ideas$ the \emph{aufheb pair} is
$\mathrm{aufheb}(a):=(\sigma(a),\nu(a))\in\Ideas\times\Ideas$.
\end{definition}

\begin{definition}[Mediation]\label{def:mediated}
For $a,b\in\Ideas$ the \emph{mediated form} is
\[
\mathrm{med}(a,b):=\sigma(a)\compose\nu(b).
\]
We read $\mathrm{med}(a,b)$ as ``the sediment of $a$ transformed by the
negation of $b$.''
\end{definition}

\begin{definition}[Atomic predicates]\label{def:aufheb-preds}
An idea $a\in\Ideas$ is
\begin{itemize}
\item \textbf{preserved} if $\sigma(a)=a$;
\item \textbf{negated} if $\nu(a)=a$;
\item \textbf{sedimented} if $\nu(a)=\unit$;
\item \textbf{sublated} if $\sigma(a)=\unit$;
\item \textbf{stable} if it is both preserved and sedimented;
\item \textbf{Hegelian} if $\nu(a)$ is sedimented, i.e.\
$\nu(\nu(a))=\unit$.
\end{itemize}
\end{definition}

\begin{definition}[Triad]\label{def:triad}
The \emph{triad} attached to $a\in\Ideas$ is the triple
\[
\mathrm{triad}(a):=(\mathrm{thesis},\mathrm{antithesis},\mathrm{synthesis})=(a,\nu(a),\sigma(a)).
\]
The triple $(a,\nu(a),\sigma(a))$ is the algebraic ghost of the
informal thesis--antithesis--synthesis schema; the recovery law
$\mathrm{synthesis}\compose\mathrm{antithesis}=\mathrm{thesis}$ holds by (A1).
\end{definition}

\begin{definition}[Iterated projectors]\label{def:iterates}
For $n\in\mathbb{N}$ define
\[
\sigma^{[n]}(a):=\underbrace{\sigma(\sigma(\cdots\sigma(a)\cdots))}_{n\text{ times}},
\qquad
\nu^{[n]}(a):=\underbrace{\nu(\nu(\cdots\nu(a)\cdots))}_{n\text{ times}},
\]
with $\sigma^{[0]}(a)=\nu^{[0]}(a)=a$.
\end{definition}

\begin{definition}[Aufheb equivalence]\label{def:aufheb-equiv}
Two ideas $a,b\in\Ideas$ are \emph{Aufheb-equivalent}, written
$a\simeq_{\text{A}} b$, if $\sigma(a)=\sigma(b)$ and
$\nu(a)=\nu(b)$.
\end{definition}

These six definitions exhaust the primitives. Everything in the rest
of the chapter is derived.

\subsection*{Closure and basic algebra}

We collect the immediate consequences of the axioms.

\begin{lemma}\label{lem:proj-fixed}
For every $a\in\Ideas$, $\sigma(a)$ is preserved and $\nu(a)$ is
negated.
\end{lemma}

\begin{proof}
By (A2), $\sigma(\sigma(a))=\sigma(a)$, so $\sigma(a)$ is a fixed
point of $\sigma$, hence preserved. By (A3), $\nu(\nu(a))=\nu(a)$,
so $\nu(a)$ is negated.
\end{proof}

\begin{lemma}\label{lem:ident-quad}
The identity $\unit$ is simultaneously preserved, negated, sedimented,
and sublated.
\end{lemma}

\begin{proof}
By (A4), $\sigma(\unit)=\unit$ and $\nu(\unit)=\unit$. The first
equation says $\unit$ is preserved (trivially) and sublated
(vacuously, since the value $\unit$ equals $\unit$). The second says
$\unit$ is negated and sedimented. All four predicates are simply
two equations $\sigma(\unit)=\unit$ and $\nu(\unit)=\unit$ read in
two ways each.
\end{proof}

\begin{lemma}\label{lem:cross-trivial}
For all $a\in\Ideas$, $\sigma(\nu(a))=\unit=\nu(\sigma(a))$.
\end{lemma}

\begin{proof}
This is (A5).
\end{proof}

\begin{lemma}\label{lem:decomp-collapse}
If $a$ is sedimented, then $\sigma(a)=a$. If $a$ is sublated, then
$\nu(a)=a$.
\end{lemma}

\begin{proof}
Suppose $\nu(a)=\unit$. By (A1), $a=\sigma(a)\compose\nu(a)
=\sigma(a)\compose\unit=\sigma(a)$ using right identity. Suppose
$\sigma(a)=\unit$. By (A1), $a=\sigma(a)\compose\nu(a)=\unit
\compose\nu(a)=\nu(a)$ using left identity.
\end{proof}

\begin{proposition}[Equivalence of preservation and sedimentation]\label{prop:pres-sed}
For all $a\in\Ideas$, $a$ is preserved if and only if $a$ is
sedimented.
\end{proposition}

\begin{proof}
($\Rightarrow$) Suppose $\sigma(a)=a$. Then
$\nu(a)=\nu(\sigma(a))=\unit$ by (A5), so $a$ is sedimented.
($\Leftarrow$) Suppose $\nu(a)=\unit$. By
\cref{lem:decomp-collapse}, $\sigma(a)=a$, so $a$ is preserved.
\end{proof}

\begin{proposition}[Equivalence of negation and sublation]\label{prop:neg-sub}
For all $a\in\Ideas$, $a$ is negated if and only if $a$ is sublated.
\end{proposition}

\begin{proof}
Symmetric to \cref{prop:pres-sed}: substitute (A5)'s second equation
$\sigma(\nu(a))=\unit$ and the left-identity collapse from
\cref{lem:decomp-collapse}.
\end{proof}

\section{The triadic decomposition theorem}\label{sec:triadic-decomp}

The five axioms package into a single structural statement. We have
already seen the individual content of each axiom in
\S\ref{sec:aufheb-defs}; the present section bundles them, addresses
the question of uniqueness modulo the kernel of the resonance
pairing $\rho$, and draws the immediate algebraic corollaries. The
distinctive feature of the bundle, philosophically, is that it
treats Aufhebung not as a process unfolding in time but as a
\emph{simultaneous} structural fact about every idea. Hegel's
\emph{Bewegung} (movement) is replaced by a static decomposition;
the temporal connotations of the German verb \emph{aufheben} are
recovered, in our setting, only when one composes successive ideas
along a chain.

\begin{theorem}[Aufhebung Decomposition Theorem]\label{thm:3-1}
Let $(\Ideas,\compose,\unit,\sigma,\nu)$ be an Aufhebung structure.
For every $a\in\Ideas$ the following five identities hold
simultaneously:
\begin{align*}
\sigma(a)\compose\nu(a) &= a, \\
\sigma(\sigma(a)) &= \sigma(a), \\
\nu(\nu(a)) &= \nu(a), \\
\nu(\sigma(a)) &= \unit, \\
\sigma(\nu(a)) &= \unit.
\end{align*}
\end{theorem}

\begin{proof}
The five identities are precisely the axioms (A1), (A2), (A3) and
the two halves of (A5), instantiated at $a$. The bundle is a
tautological consequence of \cref{def:aufheb-struct}; the
non-trivial content of the theorem is its \emph{statement}---i.e.\
the assertion that an algebraic structure satisfying just these
five conjuncts captures the full informal Hegelian schema.
\end{proof}

The deeper content is uniqueness. The bare datum
``$a=\sigma(a)\compose\nu(a)$ for some pair $(\sigma(a),\nu(a))$'' is
satisfied by infinitely many pairs in any non-trivial monoid;
declaring $\sigma(a)=a$ and $\nu(a)=\unit$, or $\sigma(a)=\unit$ and
$\nu(a)=a$, both work pointwise. What pins the pair down is the
combination of (A2)--(A5).

\begin{theorem}[Uniqueness modulo $\ker\rho$]\label{thm:uniqueness}
Let $\rho:\Ideas\times\Ideas\to\mathbb{R}$ be the resonance pairing of
\cref{ch:foundations}, and assume the algebra carries a compatible
Aufhebung structure (compatibility meaning that $\sigma$ and $\nu$
preserve $\rho$ in each argument). Suppose $\sigma',\nu'$ is a second
Aufhebung structure on the same monoid satisfying axioms (A1)--(A5)
and the same compatibility. Then for every $a\in\Ideas$,
\[
\sigma'(a)-\sigma(a)\in\ker\rho,\qquad \nu'(a)-\nu(a)\in\ker\rho,
\]
where the difference is taken in the free $\mathbb{R}$-vector space on
$\Ideas$.
\end{theorem}

\begin{proof}[Sketch]
Compose any two decompositions of $a$ by their formal difference and
pair the result with arbitrary $b\in\Ideas$ under $\rho$. By
bilinearity and compatibility,
$\rho(\sigma'(a)-\sigma(a),b)=\rho(a,b)-\rho(a,b)=0$ for all $b$, so
$\sigma'(a)-\sigma(a)$ lies in the right kernel of $\rho$; the
$\nu$-case is symmetric. The full argument is deferred to Theorem 1.3
of \cref{ch:foundations}, which we apply pointwise.
\end{proof}

The combination of \cref{thm:3-1} (existence) and
\cref{thm:uniqueness} (uniqueness) constitutes the third of the nine
foundational theorems of the SHARED BRIEF in its sharpened, structure-
level form.

\begin{corollary}[Triadic recovery]\label{cor:triad-recovery}
For every $a\in\Ideas$,
\[
\mathrm{triad}(a).\mathrm{synthesis}\compose\mathrm{triad}(a).\mathrm{antithesis}
=\mathrm{triad}(a).\mathrm{thesis}.
\]
\end{corollary}

\begin{proof}
This is (A1) restated through the names of the components of the
triad: synthesis is $\sigma(a)$, antithesis is $\nu(a)$, thesis is
$a$.
\end{proof}

\begin{lemma}[Mediation laws]\label{lem:mediation}
The mediated form satisfies, for all $a,b,c\in\Ideas$:
\begin{enumerate}
\item $\mathrm{med}(a,a)=a$;
\item $\mathrm{med}(a,\unit)=\sigma(a)$;
\item $\mathrm{med}(\unit,b)=\nu(b)$;
\item $\mathrm{med}(a,b)\compose c=\sigma(a)\compose(\nu(b)\compose c)$.
\end{enumerate}
\end{lemma}

\begin{proof}
For (1), $\mathrm{med}(a,a)=\sigma(a)\compose\nu(a)=a$ by (A1). For (2),
$\mathrm{med}(a,\unit)=\sigma(a)\compose\nu(\unit)=\sigma(a)\compose\unit
=\sigma(a)$ by (A4) and right identity. For (3),
$\mathrm{med}(\unit,b)=\sigma(\unit)\compose\nu(b)=\unit\compose\nu(b)
=\nu(b)$ by (A4) and left identity. For (4), apply the associativity
of $\compose$.
\end{proof}

\section{Identity behavior}\label{sec:aufheb-identity}

The identity $\unit$ plays a remarkable role in the Aufhebung
structure: it is the unique idea that satisfies all four atomic
predicates simultaneously. The reason is that $\unit$ is the
algebraic neutralization of composition, and composition is the
\emph{glue} that the projectors decompose; an idea that contributes
nothing to composition has, correspondingly, nothing for the
projectors to do. The identity is \emph{the} fixed point of every
piece of structure that an Aufhebung algebra carries.

\begin{theorem}[Quadruple status of the identity]\label{thm:ident-quad}
The identity $\unit\in\Ideas$ is preserved, negated, sedimented, and
sublated. Its triad is $(\unit,\unit,\unit)$ and its aufheb pair is
$(\unit,\unit)$.
\end{theorem}

\begin{proof}
Combine (A4) with the definitions: $\sigma(\unit)=\unit$ gives
preservation and sublation; $\nu(\unit)=\unit$ gives negation and
sedimentation. The triad is $(\unit,\nu(\unit),\sigma(\unit))=
(\unit,\unit,\unit)$, and the aufheb pair is the same with the first
component dropped.
\end{proof}

\begin{proposition}[Uniqueness of the quadruple-status idea]\label{prop:ident-unique-quad}
If $a\in\Ideas$ is simultaneously preserved and sublated, then
$a=\unit$. Likewise if $a$ is negated and sedimented.
\end{proposition}

\begin{proof}
Assume $a$ is preserved and sublated: $\sigma(a)=a$ and
$\sigma(a)=\unit$. Then $a=\sigma(a)=\unit$. Assume $a$ is negated
and sedimented: $\nu(a)=a$ and $\nu(a)=\unit$. Then $a=\nu(a)=\unit$.
\end{proof}

\begin{remark}
\Cref{thm:ident-quad} and \cref{prop:ident-unique-quad} together
explain why the identity is dialectically inert: it has nothing to
preserve and nothing to negate. Hegel describes \emph{Sein} (Being)
in this register at the opening of the \emph{Wissenschaft der Logik}
\citep{Hegel1812}: the most empty determination, indistinguishable
from \emph{Nichts} (Nothing), is precisely the absolute starting
point of the dialectic. In our framework $\unit$ is the unique
``moment'' for which thesis, antithesis, and synthesis coincide.
\end{remark}

There is a complementary observation for non-identity ideas. Any
$a\neq\unit$ that is preserved is, by \cref{prop:ident-unique-quad},
necessarily not sublated; any $a\neq\unit$ that is negated is
necessarily not sedimented. The four predicates therefore impose, on
$\Ideas\setminus\{\unit\}$, a strict partition into the genuinely
preserved, the genuinely negated, and the doubly indeterminate. This
partition is what the application chapters
\citep{Lovejoy1936,Bourdieu1984} appeal to when they distinguish
``settled'' ideas (case 2 above) from ``critical'' ideas (case 3)
from ``ideas in motion'' (case 4). The framework supplies a precise
algebraic referent for each.

\begin{corollary}[Identity composition with projectors]\label{cor:ident-proj}
For all $a\in\Ideas$,
\[
\sigma(a)\compose\unit=\sigma(a),\qquad
\unit\compose\nu(a)=\nu(a).
\]
\end{corollary}

\begin{proof}
The monoid identity laws applied to the projector outputs.
\end{proof}

\begin{lemma}[Cross-projector congruence]\label{lem:cross-cong}
For all $a\in\Ideas$,
\[
\nu(\sigma(a))=\nu(\unit),\qquad \sigma(\nu(a))=\sigma(\unit).
\]
\end{lemma}

\begin{proof}
By (A5), $\nu(\sigma(a))=\unit=\nu(\unit)$ using (A4); the second
follows symmetrically.
\end{proof}

\section{Idempotence and trivialization}\label{sec:idempotence}

The two projectors are not only individually idempotent but mutually
trivializing. Two readings of this fact are worth keeping apart. The
first is geometric: $\sigma$ and $\nu$ are projectors, and like
projectors on a vector space they have complementary images. Anything
in $\mathrm{Im}(\sigma)$ is invisible to $\nu$, and vice versa, with
``invisible'' meaning ``mapped to the identity.'' The second is
dialectical: an idea cannot simultaneously be ``what survives
critique'' and ``what critique strikes out''---the only idea that
plays both roles is the trivial one, and \cref{cor:mutual-exclusion}
makes this precise.

\begin{theorem}[Trivialization of crossed compositions]\label{thm:trivialization}
The compositions $\sigma\circ\nu$ and $\nu\circ\sigma$ are both the
constant function $a\mapsto\unit$.
\end{theorem}

\begin{proof}
By (A5), $\sigma(\nu(a))=\unit$ and $\nu(\sigma(a))=\unit$ for all
$a\in\Ideas$, so each composition has constant value $\unit$.
\end{proof}

\begin{corollary}[Image is a fixed-point set]\label{cor:image-fixed}
$\mathrm{Im}(\sigma)=\{a:\sigma(a)=a\}$ and
$\mathrm{Im}(\nu)=\{a:\nu(a)=a\}$.
\end{corollary}

\begin{proof}
For $\sigma$: ($\subseteq$) if $a=\sigma(b)$ for some $b$, then
$\sigma(a)=\sigma(\sigma(b))=\sigma(b)=a$ by (A2). ($\supseteq$) if
$\sigma(a)=a$, then trivially $a\in\mathrm{Im}(\sigma)$. The argument
for $\nu$ is the same.
\end{proof}

\begin{corollary}[Mutual exclusion of nontrivial images]\label{cor:mutual-exclusion}
If $a\in\mathrm{Im}(\sigma)\cap\mathrm{Im}(\nu)$ then $a=\unit$.
\end{corollary}

\begin{proof}
By \cref{cor:image-fixed}, $\sigma(a)=a$ and $\nu(a)=a$. So $a$ is
both preserved and negated. By
\cref{prop:pres-sed,prop:neg-sub}, $a$ is both sedimented and
sublated, hence $\sigma(a)=\unit$ and $\nu(a)=\unit$. Combining,
$a=\sigma(a)=\unit$.
\end{proof}

\Cref{cor:mutual-exclusion} is structurally important: it tells us
that the preserved part of any idea and the negated part of any
\emph{other} idea live in two subsets of $\Ideas$ whose intersection
is the singleton $\{\unit\}$. The decomposition (A1) thus expresses
$a$ as a product of an element of $\mathrm{Im}(\sigma)$ and an element
of $\mathrm{Im}(\nu)$, and the only way for both to equal $a$ itself
is for $a=\unit$.

We can read \cref{cor:mutual-exclusion} as a structural separation
theorem. Every idea $a\neq\unit$ falls into one of three cases.
Either $\sigma(a)\neq a$, in which case $a$ is genuinely
``transformable'' by $\sigma$ and is not preserved; or $\nu(a)\neq
a$, in which case $a$ is not negated; or both, in which case $a$
admits a non-trivial proper decomposition into preserved and negated
parts neither of which is itself. The trichotomy refines further: by
\cref{prop:pres-sed,prop:neg-sub}, ``preserved'' is the same as
``sedimented'' and ``negated'' is the same as ``sublated.'' So the
sixteen \emph{a priori} cells of the four-predicate
cross-classification reduce to four genuine cases:
\begin{enumerate}
\item $a=\unit$: preserved, negated, sedimented, sublated.
\item $a$ preserved (equivalently sedimented) but not negated.
\item $a$ negated (equivalently sublated) but not preserved.
\item $a$ neither preserved nor negated: a generic idea.
\end{enumerate}
A complete typology of an Aufhebung structure thus reduces to
counting, in each cardinality stratum of $\Ideas$, the populations
of cases (2), (3), and (4); case (1) is always a singleton.

\begin{proposition}[Quasi-projector identities]\label{prop:quasi-proj}
For all $a,b\in\Ideas$,
\[
\sigma(\sigma(a)\compose\nu(b))\text{ and }\nu(\sigma(a)\compose\nu(b))
\]
need not in general equal $\sigma(a)$ and $\nu(b)$ respectively. The
Aufhebung structure does not require $\sigma$ and $\nu$ to be
homomorphisms.
\end{proposition}

\begin{proof}
A counterexample is built in \cref{ch:foundations}'s example of the
free Aufhebung monoid on two generators: there the projectors are
defined component-wise but the underlying composition is associative
without commutation through $\sigma$ or $\nu$. The proposition
records this as a structural caveat; we will return to it when
discussing Marx's notion of \emph{inversion} in
\S\ref{sec:aufheb-zoo}.
\end{proof}

\section{The Aufheb-equivalence relation}\label{sec:aufheb-equiv}

We now study the natural equivalence relation generated by the
projectors. Recall \cref{def:aufheb-equiv}: $a\simeq_{\text{A}}b$
iff $\sigma(a)=\sigma(b)$ and $\nu(a)=\nu(b)$. The relation captures
the intuitive notion of ``two ideas with the same dialectical
anatomy''---same preserved core, same negated remainder. One would
ordinarily expect such an equivalence to have a non-trivial
quotient: many ideas could share an anatomy while differing in
synthesis. The principal result of the section is that this
expectation fails. The Aufhebung axioms force the anatomy to
determine the idea.

\begin{proposition}[Equivalence relation]\label{prop:equiv-rel}
The relation $\simeq_{\text{A}}$ on $\Ideas$ is reflexive, symmetric,
and transitive.
\end{proposition}

\begin{proof}
Reflexivity: $\sigma(a)=\sigma(a)$ and $\nu(a)=\nu(a)$. Symmetry:
swap the equations. Transitivity: if $\sigma(a)=\sigma(b)$ and
$\sigma(b)=\sigma(c)$, then $\sigma(a)=\sigma(c)$; same for $\nu$.
\end{proof}

\begin{theorem}[Collapse of the quotient]\label{thm:quotient-collapse}
$a\simeq_{\text{A}}b$ if and only if $a=b$. Equivalently, the
quotient $\Ideas/\!\simeq_{\text{A}}$ is canonically isomorphic to
$\Ideas$.
\end{theorem}

\begin{proof}
($\Leftarrow$) Trivial. ($\Rightarrow$) Suppose $\sigma(a)=\sigma(b)$
and $\nu(a)=\nu(b)$. By (A1) applied to $a$ and to $b$,
\[
a=\sigma(a)\compose\nu(a)=\sigma(b)\compose\nu(b)=b.
\]
\end{proof}

\Cref{thm:quotient-collapse} is striking: the natural equivalence
relation that one would write down to express ``two ideas with the
same Aufhebung anatomy'' turns out to coincide with strict equality.
Equivalently, the aufheb pair function
$\mathrm{aufheb}:\Ideas\to\Ideas\times\Ideas$ is injective. The injectivity
should not be mistaken for triviality. It says that within an
Aufhebung structure satisfying (A1)--(A5) every idea is reconstructed
exactly from its preserved and negated parts; there is no
``surplus'' content unaccounted for. The Hegelian programme of
characterizing an idea by its dialectical history is, in this sense,
\emph{complete} relative to a single round of decomposition.

The collapse has a corollary that one might call \emph{algebraic
holism}: the value of $a$ is determined by the joint behavior of
$\sigma$ and $\nu$ on it, with no further appeal to the carrier
$\Ideas$ as such. Two ideas that look identical to both projectors
are identical, and this is enforced by axiom (A1) alone---the four
remaining axioms are not used in the proof of
\cref{thm:quotient-collapse}. (A1) by itself thus already entails
that the decomposition is faithful; (A2)--(A5) supply the additional
algebraic content that distinguishes Aufhebung from a generic
factorization.

\begin{corollary}[Aufheb pair injectivity]\label{cor:aufheb-injective}
The function $\mathrm{aufheb}:\Ideas\to\Ideas\times\Ideas$ is injective. Its
image is contained in the subset
\[
\{(s,n)\in\Ideas\times\Ideas : \sigma(s)=s,\ \nu(n)=n,\
\nu(s)=\unit,\ \sigma(n)=\unit\}.
\]
\end{corollary}

\begin{proof}
Injectivity is \cref{thm:quotient-collapse}. The image constraint
follows from \cref{lem:proj-fixed,cor:mutual-exclusion}: $s=\sigma(a)$
is preserved hence sedimented, and $n=\nu(a)$ is negated hence
sublated.
\end{proof}

\begin{remark}
The image constraint of \cref{cor:aufheb-injective} is in general
\emph{strict}: not every pair of (preserved, negated) elements arises
as an aufheb pair, because composition $s\compose n$ may equal some
$a'\neq a$ for the unique $a$ producing them. This is the formal
shadow of \emph{aliasing} in the literary-canon corollary: two
distinct ``readings'' may share the same canonical residue and the
same labile transformation but, through their distinct compositions,
still yield distinct texts.
\end{remark}

\section{Iterated preservation and iterated negation}\label{sec:aufheb-iter}

The idempotence axioms (A2) and (A3) imply a very strong
\emph{convergence} property for iteration. Before stating it we set
up the iteration formally and check the trivial cases. The
iteration data of \cref{def:iterates} is what one would expect: a
zeroth iterate that returns the input unchanged, and a successor
iterate that applies one more projector. The technical work is to
show that the orbit collapses after one step, and that this
collapse is uniform in $n$.

\begin{lemma}[Single-step stabilization of iterates]\label{lem:single-step}
For all $n\in\mathbb{N}$ and all $a\in\Ideas$,
\[
\sigma^{[n+2]}(a)=\sigma^{[n+1]}(a),\qquad
\nu^{[n+2]}(a)=\nu^{[n+1]}(a).
\]
\end{lemma}

\begin{proof}
Induction on $n$. Base $n=0$: $\sigma^{[2]}(a)=\sigma(\sigma(a))=
\sigma(a)=\sigma^{[1]}(a)$ by (A2); same for $\nu$. Induction step:
$\sigma^{[n+3]}(a)=\sigma(\sigma^{[n+2]}(a))=\sigma(\sigma^{[n+1]}(a))
=\sigma^{[n+2]}(a)$ by the induction hypothesis applied to one less.
\end{proof}

\begin{theorem}[Closed form for iterates]\label{thm:closed-form}
For all $n\geq 1$ and all $a\in\Ideas$,
\[
\sigma^{[n]}(a)=\sigma(a),\qquad \nu^{[n]}(a)=\nu(a).
\]
\end{theorem}

\begin{proof}
Induction on $n$. Base $n=1$ is definitional. Induction step: assume
$\sigma^{[n]}(a)=\sigma(a)$ for some $n\geq 1$. By
\cref{lem:single-step} (with parameter $n-1$),
$\sigma^{[n+1]}(a)=\sigma^{[n]}(a)=\sigma(a)$. Symmetrically for
$\nu$.
\end{proof}

\begin{theorem}[Idempotent Stabilization Theorem]\label{thm:3-2}
For all $a\in\Ideas$ the following four assertions hold
simultaneously:
\begin{enumerate}
\item For every $n\geq 1$, $\sigma^{[n]}(a)=\sigma(a)$.
\item For every $n\geq 1$, $\nu^{[n]}(a)=\nu(a)$.
\item For all $m,n\geq 1$, $\sigma^{[m]}(a)=\sigma^{[n]}(a)$.
\item For all $m,n\geq 1$, $\nu^{[m]}(a)=\nu^{[n]}(a)$.
\end{enumerate}
\end{theorem}

\begin{proof}
Items 1 and 2 are \cref{thm:closed-form}. For item 3, both sides
equal $\sigma(a)$ by item 1; for item 4, both equal $\nu(a)$ by
item 2.
\end{proof}

\Cref{thm:3-2} is the formal expression of the idea, recurrent in
the dialectical tradition since \citet{Adorno1966}, that ``the
negation of the negation'' adds nothing beyond the first negation.
We obtain the conclusion in stronger form: \emph{any} positive number
of repeated preservations (or negations) collapses to the first.
There is no progress past the first sublation.

A standard objection at this point asks whether the idempotence
axioms are too strong. In the informal Hegelian text, sublation is
not a one-shot operation: each Aufhebung produces a new idea that is
itself in turn sublatable, and the dialectic is inexhaustible. We
have already seen one answer to this objection (above): iteration of
the dialectic is modeled \emph{externally}, by composing successive
ideas, not internally by iterating $\sigma$. There is also a second
answer, properly formal. Suppose we drop axiom (A2) and ask only
that $\sigma$ act as a self-map. The orbit of $a$ under $\sigma$ is
then a sequence $a,\sigma(a),\sigma^{[2]}(a),\sigma^{[3]}(a),\ldots$
of, in general, distinct elements. To say that the dialectic
\emph{converges} is to say that the orbit eventually stabilizes;
\cref{thm:3-2} gives the strongest possible form of convergence:
\emph{after one step}. Weaker forms (period $k$ stabilization,
asymptotic convergence in some metric) are conceivable but require
additional structure on $\Ideas$. The five-axiom version captures the
maximally rigid case; the deformed cases of Adorno and Derrida are
recovered by relaxing (A2) and (A3) respectively.

\begin{proposition}[Mediated iteration]\label{prop:mediated-iter}
For all $a,b\in\Ideas$ and $n\geq 1$,
\[
\sigma^{[n]}(\mathrm{med}(a,b))=\sigma^{[n]}(\sigma(a)\compose\nu(b)),
\]
which in general is \emph{not} equal to $\sigma(a)\compose\unit
=\sigma(a)$ unless $\sigma$ commutes with $\compose$.
\end{proposition}

\begin{proof}
Definitional unfolding plus a remark: the failure of equality is
\cref{prop:quasi-proj}.
\end{proof}

\begin{corollary}[Constancy on fixed points]\label{cor:fixed-point-constancy}
If $a$ is preserved, then $\sigma^{[n]}(a)=a$ for all $n\geq 0$. If
$a$ is negated, then $\nu^{[n]}(a)=a$ for all $n\geq 0$.
\end{corollary}

\begin{proof}
Induction. Base: $\sigma^{[0]}(a)=a$. Step: $\sigma^{[n+1]}(a)
=\sigma(\sigma^{[n]}(a))=\sigma(a)=a$ by hypothesis. Symmetrically
for $\nu$.
\end{proof}

\begin{corollary}[Identity orbit]\label{cor:ident-orbit}
For all $n\geq 0$, $\sigma^{[n]}(\unit)=\unit$ and
$\nu^{[n]}(\unit)=\unit$.
\end{corollary}

\begin{proof}
Apply \cref{cor:fixed-point-constancy} to $\unit$, which is both
preserved and negated by \cref{thm:ident-quad}.
\end{proof}

\subsection*{The Hegelian fixed-point theorem}

Recall (\cref{def:aufheb-preds}): $a$ is \emph{Hegelian} iff $\nu(a)$
is sedimented, i.e.\ $\nu(\nu(a))=\unit$.

\begin{lemma}[Hegelian = Sedimented]\label{lem:heg-iff-sed}
$a$ is Hegelian if and only if $a$ is sedimented.
\end{lemma}

\begin{proof}
($\Rightarrow$) If $\nu(\nu(a))=\unit$, then by (A3) $\nu(a)=\nu(\nu(a))
=\unit$, so $a$ is sedimented. ($\Leftarrow$) If $\nu(a)=\unit$, then
$\nu(\nu(a))=\nu(\unit)=\unit$ by (A4), so $a$ is Hegelian.
\end{proof}

\begin{theorem}[Hegelian Fixed-Point Theorem]\label{thm:3-3}
For every $a\in\Ideas$ the following four conditions are equivalent:
\begin{enumerate}
\item $a$ is sedimented: $\nu(a)=\unit$.
\item $a$ is preserved: $\sigma(a)=a$.
\item $a$ equals its own preserved component: $a=\sigma(a)$.
\item $a$ is Hegelian: $\nu(\nu(a))=\unit$.
\end{enumerate}
\end{theorem}

\begin{proof}
(1)$\Leftrightarrow$(2) is \cref{prop:pres-sed}. (2)$\Leftrightarrow$(3)
is the definition of preserved unfolded. (1)$\Leftrightarrow$(4) is
\cref{lem:heg-iff-sed}.
\end{proof}

\begin{corollary}[Stable = Preserved = Sedimented]\label{cor:stable-equiv}
$a$ is stable if and only if $a$ is preserved if and only if $a$ is
sedimented.
\end{corollary}

\begin{proof}
Stability requires both preservation and sedimentation; by
\cref{thm:3-3} either implies the other, so a single one suffices.
\end{proof}

\begin{corollary}[Sediments are Hegelian]\label{cor:sigma-heg}
For every $a\in\Ideas$, $\sigma(a)$ is Hegelian.
\end{corollary}

\begin{proof}
By \cref{lem:proj-fixed}, $\sigma(a)$ is preserved; by
\cref{thm:3-3}, it is therefore Hegelian.
\end{proof}

\Cref{thm:3-3} is the cleanest piece of structural information the
chapter contains. It says that within the Aufhebung axioms there is
essentially one notion of ``fully arrived'' idea, expressible in
four mutually equivalent ways: as a fixed point of $\sigma$, as a
fixed point of $\nu$ at the identity (sedimentation), as a strict
equality $a=\sigma(a)$, and as the closure under double negation
(the Hegelian predicate). What this rules out is the dialectical
hope, sometimes attributed to Hegel and routinely defended by his
nineteenth-century followers, that there is a hierarchy of
``higher'' fixed points reached only by repeated sublation:
preserved-of-preserved is a non-trivial intensification of
preserved, and so on. The theorem shows the hope is structurally
empty in any closed Aufhebung structure: \emph{preserved is the
ceiling}. Iterated dialectic does not raise an idea above
$\sigma(a)$; only composition with new ideas does.

This is the formal content of what the post-Hegelian critics
\citep{Adorno1966,Derrida1967} attack and what the orthodox
defenders \citep[in his more conservative passages]{Hegel1812}
celebrate. The theorem is neutral on which party is right; it merely
says that, whatever else is true, the four characterizations are
equivalent and the iteration of $\sigma$ stabilizes immediately.
What remains is the question of whether the closed Aufhebung is
the right model for a given application, or whether the open
$\eta$-deformed version is needed.

\subsection*{Application corollaries}

We collect four corollaries that interpret the theorems in
applications-domain language; the formal restatements were verified
in Lean as `corollary\_3\_1` through `corollary\_3\_4`. These are
not merely decorative re-readings; they identify the algebraic
content that a given applied theory must commit to in order to be
expressible in the framework. A sociologist who endorses
\cref{cor:sediment-practice} is implicitly endorsing axioms (A2)
and (A4); a literary theorist who endorses
\cref{cor:intertextual} is endorsing (A2) alone; a cognitive
scientist who endorses \cref{cor:memory} is endorsing the full
\cref{prop:pres-sed}, and so on. The corollaries thus serve as
\emph{discriminators}: the user of the framework can locate which
piece of the structure their theory does and does not commit to.

\begin{corollary}[Sediment of practice]\label{cor:sediment-practice}
In the social-science reading where an idea models a cultural
practice, the sediment of any practice after one round of cultural
reproduction equals its eventual sediment:
$\sigma^{[n+1]}(a)=\sigma(a)$ for all $n\geq 0$.
\end{corollary}

\begin{proof}
Special case of \cref{thm:closed-form}.
\end{proof}

\begin{corollary}[Intertextual stability]\label{cor:intertextual}
In the humanistic reading where an idea models a literary text, the
canonical residue is idempotent: $\sigma(\sigma(a))=\sigma(a)$.
\end{corollary}

\begin{proof}
Axiom (A2).
\end{proof}

\begin{corollary}[Memory consolidation]\label{cor:memory}
In the cognitive reading where $\sigma$ models long-term consolidation
of a memory trace and $\nu$ models the labile transformation of
recall, an idea $a$ is consolidated ($\sigma(a)=a$) iff its labile
component is null ($\nu(a)=\unit$).
\end{corollary}

\begin{proof}
\Cref{prop:pres-sed}.
\end{proof}

\begin{corollary}[Emergence invariant]\label{cor:emergence}
In the emergence-theoretic reading where $\sigma(a)$ is the
conserved core of an idea across phase transitions, $\sigma(a)$ is
itself stable.
\end{corollary}

\begin{proof}
$\sigma(a)$ is preserved by \cref{lem:proj-fixed}; preserved
implies stable by \cref{cor:stable-equiv}.
\end{proof}

\section{The zoo: seven entries}\label{sec:aufheb-zoo}

The Aufhebung decomposition of this chapter has been formulated in
algebraic vacuum. We now exhibit seven historical positions that map
onto the formalism, sometimes faithfully and sometimes
illuminatingly askew.

\subsection*{Hegel: Aufhebung literally}

\textbf{Original claim.} In the \emph{Phenomenology of Spirit}
\citep{Hegel1807}, Hegel introduces \emph{Aufhebung} as the operation
by which a determinate consciousness, on encountering its own
contradiction, does not perish but transforms itself, retaining as a
\emph{moment} what it had been. The locus classicus is \S113 of the
\emph{Phenomenology}: ``\emph{Das Aufgehobene} ist zugleich ein
Aufbewahrtes,'' the sublated is simultaneously preserved. The
\emph{Wissenschaft der Logik} \citep[\S185--\S190]{Hegel1812} systematizes
this into the structural law that \emph{Sein, Nichts, Werden} (Being,
Nothing, Becoming) form the first triad: Becoming is what
\emph{aufhebt} the contradiction of Being and Nothing, preserving
both as moments and elevating to a new determination.

\textbf{Formal restatement.} The triad
$(\mathrm{thesis},\mathrm{antithesis},\mathrm{synthesis})=(a,\nu(a),\sigma(a))$ of
\cref{def:triad} is the algebraic image of Hegel's three moments,
with the recovery law
$\mathrm{synthesis}\compose\mathrm{antithesis}=\mathrm{thesis}$ (\cref{cor:triad-recovery})
encoding the determinacy of the synthesis. Sein corresponds to a
maximally preserved $a$, Nichts to its negation, and Werden to the
canonical decomposition.

\textbf{Status.} Faithful at the structural level. The chapter
proves Theorem 3.1 (\cref{thm:3-1}), which is the algebraic version
of Hegel's claim that every determinate idea is Aufhebbar. The
sharpening lies in idempotence: Hegel speaks of an unbounded
dialectical progression, but our axioms force first-step
stabilization. The Hegelian dialectic is therefore embedded into
the structure as a \emph{single} round, with longer chains modeled
externally by composition.

\textbf{Commentary.} The discrepancy between Hegel's open-ended
progression and our \cref{thm:3-2} is the chapter's principal
divergence from the source text. We read it as a successful
formalization of Brandom's \citep{Brandom1994} reading, in which the
Hegelian process is iterated by composing successive Aufhebungen
rather than by iterating the projectors themselves.  A further point
of fidelity worth marking: the \emph{Phenomenology}'s preface
\citep[\S20]{Hegel1807} insists that the True is ``the whole''
(\emph{das Ganze}), and that the whole is the result together
with its becoming. Our \cref{cor:triad-recovery} says exactly that:
the thesis is recovered \emph{from} the synthesis composed with the
antithesis. The result and the becoming are not separated.

\subsection*{Marx: dialectical inversion}

\textbf{Original claim.} \citet{Marx1867} famously inverts Hegel:
``With him [Hegel] it is standing on its head. It must be turned
right side up again, if you would discover the rational kernel within
the mystical shell.'' Marx's dialectic conserves the structural form
of class antagonism ($\sigma$) while explicitly negating the
ideological inversion that disguises it ($\nu$). The development of
capitalism is read as a sequence of partial sublations in which the
mode of production is preserved but the legal and ideological
superstructure is critiqued and transformed.

\textbf{Formal restatement.} Identify $\sigma(a)$ with the conserved
relations of production embedded in idea $a$, and $\nu(a)$ with the
critical inversion that reveals the ideological remainder. The Marxist
move is to refuse the equation
$\sigma(a)=a$: ideas are never \emph{merely} preserved; they always
carry a non-trivial $\nu(a)\neq\unit$ component that critique must
expose. \Cref{cor:mutual-exclusion} then formalizes the requirement
that critique cannot be ``pure preservation''---an idea simultaneously
in $\mathrm{Im}(\sigma)$ and $\mathrm{Im}(\nu)$ collapses to $\unit$,
the politically vacuous identity.

\textbf{Status.} Compatible. Marx's distinction between base and
superstructure is captured by the projector pair $(\sigma,\nu)$;
the critique of ideology is the explicit construction of $\nu(a)$
where the bourgeois reading would set $\nu(a)=\unit$.

\textbf{Commentary.} The Marxian inversion is structural: $\sigma$
and $\nu$ are not symmetric in the political register. Marx
privileges $\nu$ as the operation of critique, treating $\sigma$ as
the residual ideology to be exposed. The formalism is symmetric in
its axioms but asymmetric in its applications, exactly the role Marx
assigns to dialectics versus political economy.  A subtler point
concerns the famous footnote of \citet[postface to the second
edition]{Marx1867}, where Marx distinguishes the dialectical
\emph{method} from its Hegelian \emph{form}: the method is rational,
the form is mystical. In our terms, Marx accepts axioms (A1)--(A5)
on the structure $(\sigma,\nu)$ but rejects the metaphysical
identification of $\sigma(a)$ with the ``inner essence'' of $a$.
The preserved component is, for Marx, a contingent historical
sediment, not an eternal eidetic core. \Cref{thm:closed-form} is
read accordingly: the rapid stabilization of iterated $\sigma$
captures the historical observation that ideological forms achieve a
rigid stability surprisingly quickly, but says nothing about the
\emph{content} of what gets stabilized.

\subsection*{Adorno: negative dialectics}

\textbf{Original claim.} \citet{Adorno1966} argues that the Hegelian
synthesis is always premature: there is no stable third moment that
genuinely reconciles thesis and antithesis. The rational task is to
preserve the negation as negation, refusing the move to a closed
$\mathrm{synthesis}$. Negative dialectics is dialectics minus the closure of
$\eta$.

\textbf{Formal restatement.} Adorno's stance translates to a refusal
of \cref{cor:triad-recovery}: the triad does not in fact ``recover''
$a$ as $\mathrm{synthesis}\compose\mathrm{antithesis}$; instead $\eta$ is
indefinitely deferred. Within our axioms, this corresponds to working
in a quotient where (A1) holds only up to a non-trivial residue
$r(a)\in\Ideas$ that is never zero. The Adornian programme is then
the formal study of the cohomological obstruction to (A1).

\textbf{Status.} Refuted under axiom (A1), illuminated under its
weakening. The strict Aufhebung structure of
\cref{def:aufheb-struct} satisfies (A1) by fiat; Adorno's position
is that this fiat is unjustified for substantive ideas. Our
foundational \cref{ch:foundations} \emph{does} introduce the
$\eta$-component as an independent piece of data, so Adorno's
position can be rigorously expressed by demanding $\eta\neq 0$
generically and treating the closed Aufhebung structure as the
limiting, idealized case.

\textbf{Commentary.} The chapter's strict treatment of $\sigma+\nu$
makes Adorno look like the algebraist of \emph{deformation theory}:
the synthesis is the always-deferred limit, never the present
quotient.  This is not the language Adorno would have used, but it
captures, we think, the structural content of his polemic against
the ``identity-thinking'' (\emph{Identit\"atsdenken}) of
classical dialectics. Adorno's central charge is that any closed
synthesis covers over a non-identity that resists subsumption; the
deferred $\eta$ of our framework is precisely the
non-identity that the closed Aufhebung suppresses. The Adornian
\emph{constellation}---a configuration of concepts that illuminates
its object without subsuming it---has a tidy formal counterpart: a
finite set $\{a_1,\ldots,a_k\}\subset\Ideas$ no two of which
$\simeq_{\text{A}}$-equivalence collapses to the same orbit, but
whose joint mediated forms $\mathrm{med}(a_i,a_j)$ produce a non-trivial
$\eta$-residue that articulates the object without identifying with
it.

\subsection*{Fauconnier and Turner: conceptual blending}

\textbf{Original claim.} \citet{FauconnierTurner2002} propose that
human thought routinely produces a \emph{blended space} from two (or
more) input mental spaces by selectively projecting structure from
each input into a new space that has emergent structure not present
in either. The blend is more than the union; it is a creative
synthesis.

\textbf{Formal restatement.} Identify the two input spaces with
$\sigma(a)$ and $\nu(b)$; the mediated form
$\mathrm{med}(a,b)=\sigma(a)\compose\nu(b)$ (\cref{def:mediated}) is the
blend space. The selective projection of structure corresponds to
applying $\sigma$ to one input (preserve essential frame structure)
and $\nu$ to the other (transform constituent details). The emergent
structure of the blend is the residue $\eta$ of the mediated
composition---i.e.\ the part of $\mathrm{med}(a,b)$ not predictable from
$(\sigma(a),\nu(b))$ pointwise.

\textbf{Status.} Faithful and predictive. \Cref{lem:mediation}
specifies the algebraic laws of the blend: it recovers the input
when both factors are equal ($\mathrm{med}(a,a)=a$), it degenerates to
the preserved component when one input is the identity, and it
satisfies a weak associativity through the projectors.

\textbf{Commentary.} Fauconnier and Turner do not formalize their
blends. The chapter's interpretation supplies a precise
specification: the blend's \emph{compositional skeleton} is fixed by
$(\sigma,\nu)$, and the \emph{emergent} content is exactly the
$\eta$-residue. This is, to our knowledge, the first algebraic
account of conceptual blending that derives the blend's existence
from the same axioms as the Hegelian dialectic.  Two predictions
follow. First, blends inherit the idempotence of the projectors:
once one has projected the input frames, further projection is idle
(\cref{thm:3-2}); this matches the cognitive observation that fully
elaborated blends do not benefit from further re-projection of their
inputs. Second, the injectivity of $\mathrm{aufheb}$
(\cref{cor:aufheb-injective}) implies that two distinct blends with
the same ``frame projection'' and same ``detail transformation''
must coincide as ideas; hence cognitive distinctness of blends with
the same projector signature requires non-trivial $\eta$, which is
the blend's emergent structure in the Fauconnier--Turner sense.

\subsection*{Koestler: bisociation}

\textbf{Original claim.} \citet{Koestler1964} characterizes the
creative act---humour, scientific discovery, art---as the
\emph{bisociation} of a single situation with two normally
incompatible matrices of thought. The punchline of a joke, the
\emph{Eureka} of an insight, and the metaphor of a poem all consist
in the sudden recognition of a non-trivial relationship between two
otherwise unrelated frames.

\textbf{Formal restatement.} A bisociation is a non-trivial $\eta$:
two ideas $a$ and $b$ whose individual contributions $(\sigma(a),
\nu(b))$ are predictable, but whose composition
$\sigma(a)\compose\nu(b)$ produces a residue not algebraically
derivable from $(\sigma(a),\nu(b))$ alone. The creative act is the
realization that for some pair $(a,b)$ the mediated form
$\mathrm{med}(a,b)$ has $\eta\neq 0$ in a structurally surprising
location.

\textbf{Status.} Compatible and sharpenable. Koestler's three
domains (humour, science, art) correspond to three regimes of the
$\eta$-component: humour requires $\eta$ to be \emph{abrupt}
(equation 1.7's discontinuity); science requires $\eta$ to be
\emph{persistent} (a stable new theorem); art requires $\eta$ to be
\emph{aesthetically resonant} (high $|\rho|$).

\textbf{Commentary.} The framework predicts what Koestler asserted
intuitively: that creativity is not in the inputs but in the
\emph{joint} of the inputs. Our \cref{cor:aufheb-injective} ensures
that the joint is uniquely determined by the inputs in any closed
Aufhebung structure; bisociation requires structures in which the
closure fails, i.e.\ structures with non-trivial $\eta$.  Koestler
makes much of the \emph{suddenness} of creative insight (the
``aha'' moment), and our formalism offers a structural reason: the
$\eta$-residue is, for given projector outputs, an algebraic constant
of the structure; recognizing it requires no temporal process and
hence presents itself instantaneously. The phenomenology of
discontinuity is the cognitive shadow of an algebraic discreteness.
Koestler's insistence that the same structure underlies humour,
science, and art is also recovered: in our framework all three are
reduced to the same operation $(a,b)\mapsto\mathrm{med}(a,b)$ followed
by inspection of the residue, with the three domains differing only
in the kind of $\eta$ they reward.

\subsection*{Deleuze: virtual problems}

\textbf{Original claim.} \citet{Deleuze1968} proposes that ideas, in
the Kantian transcendental sense, are not concepts but
\emph{problems}: structured fields of differential relations that
generate solutions without being identical with any of them. A
problem is virtual; its solutions actualize partially but never
exhaustively. The problem persists across all its solutions as their
common generative structure.

\textbf{Formal restatement.} A Deleuzean virtual problem is precisely
a \emph{Hegelian} idea in our sense: $\nu(\nu(a))=\unit$. Equivalently
(by \cref{thm:3-3}) it is a fixed point of $\sigma$, i.e.\ an idea
that has no labile component to actualize. Solutions of the problem
are then ideas $b$ such that $\sigma(b)=a$; they share $a$'s
preserved core and add their own $\nu(b)$.

\textbf{Status.} Faithful. \Cref{cor:sigma-heg} guarantees that for
every idea $a$ the preserved component $\sigma(a)$ is Hegelian, so
every idea has an associated virtual problem; the function
$\sigma:\Ideas\to\{\text{Hegelian ideas}\}$ is the Deleuzean
\emph{problematization} operator.

\textbf{Commentary.} Deleuze's claim that problems persist across
their solutions is exactly \cref{thm:3-2} (idempotent stabilization)
applied to the projector $\sigma$: once an idea has been
problematized, further problematization adds nothing. The virtual
field is closed under one application of $\sigma$.  A second
correspondence is more striking. Deleuze argues
\citep[ch.~IV]{Deleuze1968} that the actualization of a virtual
problem is a process of \emph{differenciation}---the production of
distinct solutions out of a common generative problem---and that
this process is asymmetric with respect to the inverse operation of
\emph{differentiation}, which moves from solutions back to problems.
The asymmetry of $\sigma$ and $\nu$ in our axioms (preserved on the
left, negated on the right) is, structurally, the same asymmetry: the
preserved component is the problem, the negated component is the
deviation by which a particular solution actualizes it, and the two
do not commute.

\subsection*{Derrida: différance and the deferral of $\eta$}

\textbf{Original claim.} \citet{Derrida1967} introduces
\emph{différance}---a coined word combining ``difference'' and
``deferral''---to name the play of signification that prevents any
sign from achieving final, present meaning. Meaning is always
displaced: each signifier defers to others, and full presence is
indefinitely postponed. There is no closed synthesis of a sign with
its trace.

\textbf{Formal restatement.} Derrida's position is the structural
denial of \cref{thm:quotient-collapse}: the aufheb pair
$\mathrm{aufheb}(a)=(\sigma(a),\nu(a))$ does not in general suffice to
reconstruct $a$, because the act of reconstruction
$\sigma(a)\compose\nu(a)$ itself produces a trace that is then
written into a new $a'\neq a$. Différance is the refusal of the
synthesis $\eta$ to terminate.

\textbf{Status.} Refuted under axiom (A1), reframed under its
weakening. As with Adorno, the strict Aufhebung structure forces
$\eta$ to be implicit and immediately closed; Derrida's claim is that
$\eta$ has its own dynamics. The framework accommodates this by
treating closed Aufhebung as a limiting idealization and the open
case as the generic phenomenon.

\textbf{Commentary.} Derrida and Adorno occupy structurally similar
positions in our framework: both refuse the closure of (A1). They
diverge on what to do with the residue: Adorno treats it as a
political-philosophical resource to be defended; Derrida treats it
as a textual-graphic effect to be tracked across signifying chains.
The chapter offers no verdict between them; both are formal
articulations of the open-Aufhebung regime in which (A1) holds only
up to a non-trivial $\eta$.  We add one technical observation. The
\emph{trace} of \citet{Derrida1967}, the mark of an absent term that
survives in a present sign, is naturally modeled as the difference
$a-\sigma(a)\compose\nu(a)$ in the free vector space on $\Ideas$.
Under axiom (A1) this difference vanishes; weakening (A1) to
``vanishes modulo $\ker\rho$'' (\cref{thm:uniqueness}) lets the trace
live, but only in the directions in which $\rho$ does not see it.
Différance is then exactly the kernel of $\rho$: the part of an idea
that no resonance test can detect, but that the algebra demands as
the residue of incomplete decomposition. This is a precise sense in
which Derrida is right that signification leaves a trace that
representation cannot capture.

\bigskip

The seven entries traverse two centuries of dialectical thought and
reduce them, with varying degrees of fidelity, to predicates and
operations on a single algebraic structure. The mathematics has not
displaced the philosophy; it has supplied a notation in which the
philosophical disagreements become \emph{structural}, hence
discussable. Hegel and Marx differ on the political asymmetry of
$\sigma$ versus $\nu$. Adorno and Derrida differ from Hegel on the
closure of (A1). Fauconnier--Turner and Koestler agree on the
non-triviality of $\eta$ but disagree on whether to call its
products ``blends'' or ``bisociations.'' Deleuze names the
$\sigma$-fixed points and calls them virtual problems. The next
chapter takes up the emergence cocycle $\omega$, which measures
\emph{how badly} (A1) can fail when one drops it---and so makes the
Adornian and Derridean positions algebraically precise.

We close with a methodological observation. The chapter has shown
that a five-axiom structure suffices to recover, in a recognizable
algebraic form, the principal moments of two centuries of
dialectical theorizing. This is a small piece of evidence for the
broader programmatic claim of Idea Theory, that the apparently
incommensurable vocabularies of philosophy, sociology, cognitive
science, and literary theory are unified at the level of structural
operations on a common algebraic object. The Aufhebung structure of
this chapter is one such common object; the resonance pairing of
\cref{ch:foundations} is another; the emergence cocycle of the next
chapter is a third. Each is a piece of structure that any sufficiently
substantive theory of ideas must, on this hypothesis, contain---if
only implicitly. Making the implicit explicit is the work of the
formalism. The work of the zoo is to verify that the explicit
formalism does in fact reconstruct what the original theories were
trying to say. The reader who has followed the chapter to this point
should now be in a position to recognize, in the next two chapters'
treatment of the cocycle $\omega$ and the curvature $\kappa$, the
same pattern of disagreement and agreement that has organized the
present one: a small set of axioms, a battery of elementary lemmas
deciding everything that the axioms entail, and a zoo of historical
positions whose relationship to the axioms is now structural rather
than rhetorical.
